% ============================================
% Chapter 3: System Analysis
% ============================================

\chapter{System Analysis}
\label{ch:system-analysis}

GreenLab --- Chapter 3: System Analysis





\section{Introduction to the Chapter}

This chapter is the analytical foundation on which the design and implementation of GreenLab,
described in the subsequent chapters, are built. All technical values, requirement identifiers,
database fields, and API paths quoted in this chapter are taken directly from the project's TA
Reference Document and are treated as authoritative.

Chapter 2 reviewed the literature on occupancy-driven building automation and identified a gap
that no single reviewed system closes: the combination of multimodal embedded sensing, YOLObased visual occupancy counting, a rule-based backend automation engine, a role-based
governance dashboard, immutable audit logging, and laboratory transfer management. This chapter
turns that gap into a concrete, analysable system. It answers three questions in sequence. First,
what is GreenLab, and what kind of system is it --- Section 3.2 positions GreenLab as a four-tier
cyber-physical system and explains the temporal and governance pressures that shape every
requirement that follows, and Section 3.3 immediately translates that system context into a
commercial case using the Business Model Canvas. Second, who needs GreenLab to do what, and
how do they interact with it --- Section 3.4 identifies the six stakeholder groups affected by the
system, and Section 3.5 translates their needs into a formal use case model. Third, precisely what
must the system do and how well must it do it --- Sections 3.6 and 3.7 state the 36 functional
requirements and the non-functional requirements, organised by area and quality dimension
respectively, exactly as fixed in the TA Reference Document.

The chapter closes by validating the requirement set dynamically rather than only statically. A
requirements table states what the system must do in isolation; it does not show how those
requirements compose into an end-to-end behaviour that crosses all four tiers in a single
operational scenario. Section 3.8 supplies that missing dimension by walking through the four
sequence diagrams produced for GreenLab --- the occupancy-triggered automation workflow, the
live camera streaming workflow, the manual override workflow, and the laboratory transfer request
workflow --- explaining each lifeline, each conditional branch, and each loop in terms of the
functional requirements and database tables it realises. Section 3.9 then summarises the chapter and
states explicitly how its outputs feed into the design chapter that follows.
1
GreenLab --- Chapter 3: System Analysis


Two conventions are used consistently throughout this chapter and the rest of the project book.
First, the four tiers are always referred to by their fixed names --- the hardware layer, the Python AI
service, the ASP.NET Core backend, and the Angular dashboard --- rather than informal synonyms
such as ``the server'' or ``the frontend'', in line with the project's terminology standard. Second,
every functional and non-functional requirement is cited by its fixed identifier (for example FRLC-03 or NFR-SEC-01) so that it can be traced unambiguously into the design, implementation,
and testing chapters.




2
GreenLab --- Chapter 3: System Analysis



\section{System Context}

GreenLab is analysed as a cyber-physical system that bridges the physical environment of a
university IT laboratory with software services that automate energy management and provide
governance tools for authorised users. The system is organised into four tiers, each owning a
distinct responsibility and exposing a well-defined interface to its neighbours, as summarised in
Table 3.1.

Table 3.1: GreenLab four-tier architecture and tier responsibilities.

Tier Exact Term Technology Responsibility

Tier 1 the hardware layer ESP32 + Sensors + Relay Sensing, relay control,
entry/exit counting

Tier 2 the Python AI service Python + YOLOv11n Visual occupancy detection,
person counting

Tier 3 the ASP.NET Core backend ASP.NET Core 8 + SQL Automation engine, API,
Server RBAC, audit logging

Tier 4 the Angular dashboard Angular SPA Real-time monitoring,
override, admin governance




Figure 3.1: GreenLab four-tier system architecture and data flow.

3
GreenLab --- Chapter 3: System Analysis


The hardware layer, built around the ESP32 microcontroller, observes laboratory conditions
through three connected sensors --- a DHT11 temperature and humidity sensor, an LDR lightdependent resistor, and an HC-SR04 ultrasonic distance sensor --- and acts on the physical
environment through a four-channel relay module that switches two lighting circuits and two
ventilation fans. The Python AI service runs YOLOv8n object detection on frames pulled from an
IP camera, providing a person count that complements the hardware tier's entry/exit counting. The
ASP.NET Core backend applies the automation rules captured in the decision matrix and fan
activation table, persists every state change and event to SQL Server, and exposes the 21
authenticated REST endpoints listed in Section 3.6 to both the frontend and the hardware devices.
The Angular dashboard presents real-time laboratory monitoring, manual override controls,
scheduling, transfer-request forms, and administrative tools to instructors and administrators.

This separation of concerns is intentional and is exploited throughout the project. The hardware
layer can adopt a new sensor without the backend's rule engine changing; the Python AI service
could be replaced with a different detection model without the Angular dashboard ever noticing;
the ASP.NET Core backend can add new endpoints without any firmware change; and the Angular
dashboard can be redesigned without touching server-side code. Feature dependencies in Table 3.2
make this layering explicit: every frontend feature depends on authentication (BE-03); the
automation engine (BE-01) depends on live sensor data (HW-01) and on both occupancy signals
(AI-01, AI-02); relay-based appliance control (HW-02) is only ever driven by the automation
engine or by an override (BE-01, BE-02); and manual override itself (BE-02) depends on
authentication and produces audit log entries (BE-03, BE-04).

Table 3.2: Feature dependency chain across the four-tier architecture.

Feature(s) Depends On Reason

FE-01, FE-02, FE-03, FE-04 BE-03 All frontend features require
authentication (JWT)

BE-01 HW-01, AI-01, AI-02 Automation engine needs live sensor +
occupancy data

HW-02 BE-01, BE-02 Relay control driven by automation
engine or override




4
GreenLab --- Chapter 3: System Analysis


Feature(s) Depends On Reason

BE-02 BE-03, BE-04 Override requires auth and generates
audit log entries


\subsection{Temporal Dimension}

Unlike residential or commercial buildings, where occupancy patterns are comparatively stable, a
university IT laboratory experiences highly variable occupancy that changes by semester, by
session, and sometimes within a single day: a taught session from 8:00 to 10:00, two empty hours,
an open-lab period from 12:00 to 14:00, and an examination from 14:00 to 16:00 are all plausible
within one room on one day. Each usage type implies different lighting and ventilation needs --- a
projector-led taught session benefits from dimmed lighting, while an open-lab period with active
workstations needs full lighting --- so the system must combine real-time, sensor-driven
automation with schedule-aware context rather than a single static rule set.

\subsection{Governance Dimension}

GreenLab is deliberately not a simple sensor-to-relay automation system. University laboratory
operations are governed by institutional policy that defines who may control which appliances,
how a laboratory transfer request is reviewed, what must be auditable, and how a fault is escalated.
These governance requirements are why the functional requirement set in Section 3.6 includes a
dedicated authentication and access control area (FR-AC), a manual override area with explicit
precedence rules between instructor and administrator (FR-MO), and an INSERT-only AuditLog
table that makes every state-changing action attributable and tamper-evident.

\subsection{Resilience and Safe Defaults}

Because GreenLab is cyber-physical, it must keep behaving sensibly when the network is
unreliable, when sensor readings are noisy, or when the Python AI service is temporarily
unreachable. The design therefore commits to a safe default at every tier: the ESP32 holds its last
known relay state without resetting when it loses contact with the backend (NFR-REL-01); the
backend continues serving cached sensor data and honouring active overrides if the AI service is
offline; and when any fault is active, the automation engine holds Relay 1 (lighting) ON and leaves
Relays 2--4 at their last commanded state rather than guessing (FR-FH-02). Grace periods --- the 3second HC-SR04 debounce (FR-EE-02) and the 15/30-second sensor and camera timeout


5
GreenLab --- Chapter 3: System Analysis


thresholds (FR-FH-01) --- prevent transient sensor noise or brief connectivity gaps from being
mistaken for genuine occupancy changes or hardware failure.

Read together, the system context is a set of four directional data flows governed by a single
authorisation point at the backend: sensor data flows upward from the hardware layer; occupancy
data flows upward from the Python AI service; control commands flow downward from the
Angular dashboard, through the backend, to the hardware layer; and audit records flow from every
tier into SQL Server. Sections 3.6 to 3.8 formalise this context first as testable requirements and
then as concrete sequence diagrams.




6
GreenLab --- Chapter 3: System Analysis



\section{Business Model}

Section 3.2 established GreenLab as a technical system: its four-tier architecture and the temporal
and governance pressures that shape it. A graduation project intended for genuine deployment
beyond the university, however, must also be positioned as a viable proposition rather than a
working prototype alone. This section supplies that commercial dimension using the Business
Model Canvas (BMC), a strategic framework that maps how an organisation creates, delivers, and
captures value across nine interlocking blocks. Figure 3.2 presents the completed canvas for
GreenLab; the discussion that follows explains each block in turn and traces it back to the four-tier
architecture of Section 3.2 and to the stakeholder, requirement, and sequence-diagram analysis
presented later in this chapter, so that the business case is grounded in the same requirement set as
the rest of the chapter rather than treated as a separate exercise.




Figure 3.2: GreenLab Business Model Canvas.

\subsection{Value Proposition and Customer Focus}

The value propositions in the canvas restate, in commercial language, exactly what the requirement
areas fixed later in this chapter (Section 3.6) express in technical language. Reduced energy
consumption, lower operating costs, and the reduction of electricity waste are the direct
commercial expression of the Lighting Control and Ventilation Control areas (FR-LC, FR-VC),
which switch relays off whenever occupancy is Empty rather than leaving appliances running by

7
GreenLab --- Chapter 3: System Analysis


default. Automating lighting and ventilation and monitoring laboratory occupancy in real time map
directly onto the automation engine and the Monitoring and Dashboard area (FR-MD): a laboratory
operator is not asked to trust an opaque automation engine, but can see its occupancy and relay
state at any moment, consistent with the refresh behaviour fixed by FR-MD-01. Taken together,
these value propositions are framed in the canvas as supporting a single overarching goal ---
sustainable campus operations --- which positions GreenLab as a campus-wide energy-efficiency
initiative rather than only a per-laboratory monitoring tool.

The customer segments named in the canvas are a refinement of the six stakeholder groups
identified in Section 3.4, now split explicitly into a primary and a secondary segment. Universities
and educational institutions form the primary segment: the buying organisations that would
authorise a GreenLab deployment and its budget, corresponding to the Faculty Management row of
the stakeholder table. IT departments, laboratory staff, facilities and maintenance teams, and
university administrators form the secondary segment: the operational users who interact with the
system day to day or who depend on its outputs, corresponding respectively to the IT Department,
Laboratory Staff, and Administrator rows of the same table. Separating the two makes explicit that
GreenLab is sold to an institution as a single purchasing decision, even though several distinct roles
within that institution go on to operate it. The customer relationship block completes this picture on
the service side: technical support, installation service, scheduled maintenance, software updates,
and training sessions are the ongoing obligations a vendor of GreenLab would owe its customers,
and the training-session commitment in particular is what keeps the low-friction override path of
NFR-USA-01 credible for an instructor who is a subject-matter expert rather than an IT specialist.

\subsection{Key Partners, Activities, and Resources}

The key partners block identifies the organisations a real GreenLab deployment depends on outside
the development team itself: the university administration and its facilities and maintenance
department as the institutional sponsor and physical site owner, the IT department as the network
and security gatekeeper referenced throughout Section 3.2.2's governance discussion, IoT
hardware vendors supplying the ESP32 boards, sensors, and relay modules, and a cloud or server
hosting provider for the ASP.NET Core backend and SQL Server database. Rather than assuming a
dedicated camera infrastructure provider, the canvas treats the IP cameras already installed across
university laboratories as an existing asset that the Python AI service reuses, which lowers the
barrier to a pilot deployment considerably. Each of these partners corresponds to a tier or a cross-
8
GreenLab --- Chapter 3: System Analysis


cutting concern already fixed in Table 3.1: IoT hardware vendors supply Tier 1, the existing
university camera estate supplies the visual input consumed by Tier 2, and the hosting provider
supplies the infrastructure Tier 3 runs on.

The key activities --- energy management, occupancy counting, lighting and ventilation
automation, and real-time system monitoring --- are the commercial names for the automatic use
cases detailed later in this chapter (UC-09 through UC-13): they are the activities the system
performs continuously regardless of who is watching, and they are what a customer is ultimately
paying to have run reliably. The key resources block separates the assets that make those activities
possible into three categories that mirror the architecture directly: physical resources (cameras and
IoT devices) correspond to Tier 1 and the AI-processing camera input; technology resources (the
YOLO model and the SQL Server database) correspond to Tier 2 and the persistence layer
described in Section 3.2; and human resources (software developers) are the team responsible for
maintaining the backend's decision matrix, fan activation table, and Angular dashboard as the
system evolves beyond the pilot deployment analysed in this chapter.

\subsection{Channels}

The channels block sets out how GreenLab would reach and be adopted by the customer segments
identified above: direct sales to faculty management, university partnerships that allow a pilot
deployment in one laboratory to expand across a faculty's full estate (the scenario NFR-SCAL-01
is designed to support), participation in technology competitions that provide visibility among IoT
solution providers, and relationships with existing IoT solution providers who could integrate
GreenLab's automation engine alongside their own hardware offerings. These channels are
consistent with the project's current status as a graduation project rather than a commercial
product: the university partnership and competition channels are the most immediately actionable
routes, while direct sales and third-party integration represent the longer-term path once the fourtier architecture has been validated at the scale Section 3.7.6 targets.

\subsection{Cost Structure and Revenue Streams}

The cost structure groups GreenLab's expenses into three areas that again align with the
architecture: hardware costs (the ESP32 boards, sensors, relays, and cameras that make up Tier 1
and the camera input to Tier 2), software development costs (backend development, dashboard
development, and AI maintenance and integration, spanning Tier 2 through Tier 4), and cloud and

9
GreenLab --- Chapter 3: System Analysis


infrastructure costs (database hosting for the SQL Server instance described in Section 3.2). Set
against these costs, the revenue streams propose five ways GreenLab could recover its costs and
generate a margin: one-time installation fees for initial system setup; annual maintenance contracts
that fund ongoing support; a recurring software subscription that covers access to the monitoring
dashboard and future updates; paid custom solutions for client-specific requirements; and
integration services for connecting GreenLab to a customer's existing cameras and IoT devices
rather than requiring a full hardware replacement. The subscription and maintenance-contract
streams in particular depend on the reliability targets of Section 3.7.2 --- a backend that is
unavailable during operating hours, or an ESP32 that resets its relay state on every network
interruption, would undermine the recurring-revenue case as directly as it would undermine the
technical case for the system.

Table 3.3 closes this section by tracing each canvas block back to the chapter section or
requirement area it draws on, preserving the same traceability discipline applied to the use cases
and sequence diagrams later in this chapter.

Table 3.3: Business Model Canvas blocks traced to their chapter basis.

BMC Block Chapter Basis

Value Propositions FR-LC, FR-VC, FR-OD, FR-MO, FR-MD, NFR-SCAL

Customer Segments Table 3.4 (Stakeholder Analysis); Faculty Management
(primary), IT Department / Laboratory Staff / Administrator
(secondary)

Customer Relationships NFR-USA-01, NFR-MAIN-01

Key Partners Table 3.1 (four-tier architecture); Section 3.2.2

Key Activities UC-09--UC-13 (Section 3.5.3)

Key Resources Section 3.2 (tier responsibilities)

Channels Section 3.7.6 (Scalability)

Cost Structure Table 3.1 tier technologies

Revenue Streams Section 3.7.2 (Reliability)



10
GreenLab --- Chapter 3: System Analysis


Read alongside the rest of this chapter, the canvas in Figure 3.2 shows that GreenLab's
requirement set was not defined in a commercial vacuum: the same automation behaviour,
governance controls, and reliability targets that satisfy the functional and non-functional
requirements set out later in this chapter are also the features and guarantees a paying customer
would evaluate before adopting the system beyond a single pilot laboratory.




11
GreenLab --- Chapter 3: System Analysis



\section{Stakeholder Analysis}

Six stakeholder groups have a direct or indirect interest in GreenLab's operation. Identifying them
precisely is what keeps the requirements in Sections 3.6 and 3.7 grounded in real operational need
rather than hypothetical scenarios.

Table 3.4: Stakeholder analysis summary.

Stakeholder Role Primary Concern System Features

Instructor Conducts sessions, monitors Automatic control must not Real-time dashboard,
assigned labs, overrides interfere with teaching manual override with timed
automation when teaching expiry, transfer request form
needs it, requests transfers

Administrator Manages users, schedules, Operational oversight, User management, schedule
transfers, and configuration; policy compliance, management, transfer
reviews audit logs and faults accountability approval, audit log, fault
acknowledgement

Laboratory Staff Monitors hardware integrity, Hardware faults detected Fault alert dashboard, sensor
responds to fault alerts, and resolved before they status display, relay
maintains equipment compromise safety feedback

Student Occupies the laboratory; an A comfortable environment Indirect benefit through
indirect beneficiary without needing to interact occupancy-driven lighting
with the system and ventilation

IT Department Maintains network, server, Reliable infrastructure, HTTPS, JWT, service
and security infrastructure security compliance, accounts, firmware update
updateable firmware support, network
configuration

Faculty Management Reviews energy trends and Energy cost reduction, Energy reports, audit log
policy compliance; environmental review, configuration
authorises budget responsibility, auditability oversight

The instructor and administrator are the two roles that interact directly with the Angular dashboard
and therefore drive most of the use case model in Section 3.5; laboratory staff, students, the IT
department, and faculty management are served indirectly through the features that the instructor
and administrator operate on their behalf --- fault visibility for laboratory staff, ambient comfort for

12
GreenLab --- Chapter 3: System Analysis


students, HTTPS/JWT compliance for the IT department, and audit/reporting access for faculty
management.




13
GreenLab --- Chapter 3: System Analysis



\section{Use Case Analysis}

The use case diagram for the GreenLab Intelligent Energy Management System (Figure 3.3)
captures every interaction identified in the stakeholder analysis above and groups them around four
actors: the two human roles, Instructor and Administrator, and two automated actors, the Hardware
System (Automatic) and the AI Processing System. Separating human-triggered use cases from
automatically-triggered ones makes explicit which behaviours require a logged-in user and which
run continuously regardless of who is watching --- a distinction that matters because, as Section
\subsection{explained, GreenLab's governance model hinges on knowing exactly which actions are}

attributable to a person.




Figure 3.3: GreenLab use case diagram showing actors and system interactions.

\subsection{Actors}

\begin{itemize}
  \item Instructor --- a human actor who logs in to monitor their assigned laboratory, issue manual
overrides, and submit transfer requests.
\end{itemize}

\begin{itemize}
  \item Administrator --- a human actor with faculty-wide visibility who manages user accounts,
laboratory schedules, transfer approvals, and system logs.
\end{itemize}

\begin{itemize}
  \item Hardware System (Automatic) --- the ESP32 and its sensors and relay module, acting
without a human in the loop to apply occupancy-based lighting and temperature-based
ventilation control.
14
GreenLab --- Chapter 3: System Analysis
\end{itemize}


\begin{itemize}
  \item AI Processing System --- the Python YOLO-based occupancy detection service, acting as
an automated actor that supplies person-count data into the same automatic control loop.
\end{itemize}

\subsection{Human-Triggered Use Cases}

Both human actors share ``Login to Dashboard'' as a common entry point. The diagram models this
with <<include>> relationships: ``Manual Override Control Lights and Fans'', ``Submit Laboratory
Transfer Request'', and ``Manage Laboratory Schedule'' each include ``Login to Dashboard'',
meaning none of them can be performed without a valid, authenticated session --- the direct use
case expression of FR-AC-02 (every endpoint except login requires a valid JWT).

\begin{itemize}
  \item Login to Dashboard --- validates credentials and returns a role-bearing JWT (realised by
FR-AC-01).
\end{itemize}

\begin{itemize}
  \item View Real-Time Monitoring Data --- the instructor or administrator views live LDR,
temperature, humidity, occupancy, and relay state for a laboratory (FR-MD-01, FR-MD-
03).
\end{itemize}

\begin{itemize}
  \item Manual Override Control Lights and Fans --- issues a timed override for a specific relay
(FR-MO-01, FR-MO-02); this is the use case realised in detail by the sequence diagram in
Section 3.8.3.
\end{itemize}

\begin{itemize}
  \item Submit Laboratory Transfer Request --- an instructor requests that a session move to a
different laboratory, creating a row in TransferRequests via POST /transfer/request.
\end{itemize}

\begin{itemize}
  \item Manage Laboratory Schedule --- an administrator updates a laboratory's session schedule
via PUT /labs/{id}/schedule.
\end{itemize}

\begin{itemize}
  \item Approve/Reject      Transfer   Requests        ---   an   administrator    reviews     a    pending
TransferRequests row and resolves it via PUT /transfer/{id}/approve or PUT
/transfer/{id}/reject; the diagram models the approval decision as an <<extend>> of the base
review action, since rejection is an alternative outcome of the same review rather than a
separate trigger.
\end{itemize}

\begin{itemize}
  \item Manage User Accounts --- an administrator creates or updates Users rows via POST /users
and PUT /users/{id}.
\end{itemize}



15
GreenLab --- Chapter 3: System Analysis


\begin{itemize}
  \item View System Logs and Reports --- an administrator inspects the INSERT-only AuditLog
table and active FaultEvents for oversight and incident investigation.
\end{itemize}

\subsection{Automatic Use Cases}

The lower half of the diagram belongs to the two automated actors and is connected back into
``Login to Dashboard'' by dashed <<include>> arrows. This is not a contradiction of the ``no human
in the loop'' description above --- it represents the fact that the automatic control loop's output
(current relay state, current occupancy classification) is what the dashboard displays once a user
does log in, so the automatic use cases feed the same dashboard rather than triggering it.

\begin{itemize}
  \item Automatic Occupancy-Based Lighting Control --- the decision procedure behind FR-LC-
01 and FR-LC-03, evaluated on every 5-second sensor cycle by the Hardware System actor
in conjunction with the backend's decision matrix.
\end{itemize}

\begin{itemize}
  \item Automatic Temperature-Based Ventilation Control --- the decision procedure behind
FR-VC-01 and FR-VC-02, giving the Hot temperature zone priority over occupancy when
deciding fan state.
\end{itemize}

\begin{itemize}
  \item Continuous Monitoring --- the unending 5-second sensor polling cycle (FR-SI of Section
\end{itemize}
\subsection{onward) that keeps SensorReadings current and is the precondition for every other}

automatic use case.

\begin{itemize}
  \item End of Session Reset --- the behaviour invoked when occupancy returns to Empty with no
active override, returning both light relays and, where the temperature zone permits, both
fan relays to OFF (FR-LC-03, FR-VC-03).
\end{itemize}

\begin{itemize}
  \item Perform Occupancy Detection (YOLO-based) --- owned by the AI Processing System
actor; this is FR-OD-01 through FR-OD-04, and it is the use case explained step by step in
the sequence diagram of Section 3.8.1.
\end{itemize}




16
GreenLab --- Chapter 3: System Analysis


\subsection{Use Case Summary Table}

Table 3.5 and Table 3.6 together trace every use case identified above to the functional
requirements it realises, split across the two tables below for readability.

Table 3.5: Use case descriptions and traceability to functional requirements (UC-01 to UC-06).

Traced
ID Use Case Actor(s) Description Requirements

UC-01 Login to Dashboard Instructor, Authenticate and FR-AC-01, FR-AC-
Administrator receive a role-scoped 02
JWT.

UC-02 View Real-Time Instructor, Display live sensor, FR-MD-01--FR-MD-
Monitoring Data Administrator occupancy, and relay 04
state for one or all
labs.

UC-03 Manual Override Instructor, Issue a timed override FR-MO-01--FR-MO-
Control Lights and Administrator on a specific relay, 04, FR-LC-02, FR-
Fans superseding VC-04
automation.

UC-04 Submit Laboratory Instructor Request a session TransferRequests
Transfer Request move to a different table; Endpoint 13
laboratory.

UC-05 Manage Laboratory Administrator Create or update a Endpoint 12
Schedule laboratory's session
schedule.

UC-06 Approve/Reject Administrator Review and resolve a Endpoints 14--16
Transfer Requests pending transfer
request.




17
GreenLab --- Chapter 3: System Analysis


Table 3.6: Use case descriptions and traceability to functional requirements (UC-07 to UC-13).

Traced
ID Use Case Actor(s) Description Requirements

UC-07 Manage User Administrator Create or update user Endpoints 17--19, FR-
Accounts accounts and role AC-04
assignment.

UC-08 View System Logs Administrator Inspect the audit trail FR-FH-03, FR-FH-
and Reports and active fault 04, AuditLog
events.

UC-09 Automatic Hardware System Apply the lighting FR-LC-01, FR-LC-
Occupancy-Based (Automatic) decision matrix on 03, FR-LC-04
Lighting Control every sensor cycle.

UC-10 Automatic Hardware System Apply the fan FR-VC-01--FR-VC-
Temperature-Based (Automatic) activation table on 04
Ventilation Control every sensor cycle.

UC-11 Continuous Hardware System Poll all three sensors FR-SI-01--FR-SI-04,
Monitoring (Automatic) every 5 seconds. FR-EE-01--FR-EE-04

UC-12 End of Session Reset Hardware System Return relays to OFF FR-LC-03, FR-VC-03
(Automatic) once occupancy is
Empty and no
override is active.

UC-13 Perform Occupancy AI Processing System Run YOLOv8n FR-OD-01--FR-OD-
Detection (YOLO- inference and report a 04
based) classified person
count.




18
GreenLab --- Chapter 3: System Analysis



\section{Functional Requirements}

The functional requirements are organised into the nine areas defined in the TA Reference
Document, with four requirements per area, for a total of thirty-six. Requirement identifiers follow
the pattern FR-[AREA]-[NUMBER]; priority is marked P1 (must have) or P2 (should have). These
nine areas correspond directly to the use cases of Section 3.5 and to the sequence diagrams
analysed in Section 3.8: System Initialisation feeds Continuous Monitoring (UC-11); Entry/Exit
Detection and Occupancy Detection together feed the occupancy-triggered automation workflow of
Section 3.8.1; Lighting and Ventilation Control are the two automatic use cases UC-09 and UC-10;
Monitoring and Dashboard underlies UC-02 and the camera streaming workflow of Section 3.8.2;
Manual Override is UC-03 and the workflow of Section 3.8.3; Fault Handling and Safety underlies
UC-08; and Authentication and Access Control underlies UC-01 and gates every other use case in
the model.

\subsection{Area: System Initialisation and Startup}

These requirements govern the ESP32 firmware's behaviour during power-up, before the device
enters its normal sensing-and-control loop. They establish the conditions under which the hardware
tier is considered ready to participate in the wider GreenLab data flow, and they define the fallback
behaviour when network connectivity cannot be established immediately.

Table 3.7: Area: System Initialisation and Startup.

FR-ID Requirement Priority

FR-SI-01 The ESP32 shall verify all three sensor P1
interfaces on power-up before entering
the normal polling loop; if any sensor
fails, log fault and proceed.

FR-SI-02 The ESP32 shall attempt Wi-Fi P1
connection with up to 3 retries on
startup; if all fail, enter Degraded
Mode with all relays holding last state.




19
GreenLab --- Chapter 3: System Analysis


FR-ID Requirement Priority

FR-SI-03 On first successful Wi-Fi connection P1
after startup, the ESP32 shall retrieve
the current active relay command from
GET /relay/state/{labId}.

FR-SI-04 The backend shall reject any sensor P1
POST from an ESP32 presenting an
invalid or expired Service-role JWT,
returning HTTP 401.


\subsection{Area: Entry and Exit Detection}

Entry and exit detection is the hardware tier's first line of occupancy awareness, implemented
through the HC-SR04 ultrasonic sensor. These requirements define how a physical presence event
at the laboratory doorway is converted into a debounced, countable signal that feeds the occupancy
counter referenced throughout the automation engine.

Table 3.8: Area: Entry and Exit Detection.

FR-ID Requirement Priority

FR-EE-01 The ESP32 shall classify an entry P1
event when an HC-SR04 reading drops
below 200 cm and increment the local
occupancy counter.

FR-EE-02 The ESP32 shall enforce a minimum P1
3-second debounce interval between
consecutive HC-SR04 entry or exit
events.

FR-EE-03 The ESP32 shall include the current P2
occupancy counter value in every
sensor POST payload sent to
/sensor/readings.




20
GreenLab --- Chapter 3: System Analysis


FR-ID Requirement Priority

FR-EE-04 If the HC-SR04 returns no echo within P2
30 ms (out of effective range), the
ESP32 shall record the reading as null
in SensorReadings.


\subsection{Area: Occupancy Detection and Classification}

Where entry/exit counting gives a coarse occupancy signal, the Python AI service provides a
richer, vision-based estimate using YOLOv8n person detection. These requirements define the
inference cycle, the classification thresholds that translate a raw person count into an operational
category, and the reconciliation rule applied when the two occupancy signals disagree.

Table 3.9: Area: Occupancy Detection and Classification.

FR-ID Requirement Priority

FR-OD-01 The Python AI service shall run P1
YOLOv8n inference every 5 seconds
per lab, filter detections to class 0 with
confidence $\ge$ 0.6.

FR-OD-02 The ASP.NET Core backend shall map P1
the received person count to exactly
one of four classifications: Empty (0),
Low (1--5), Moderate (6--15), High
(>15).

FR-OD-03 If the IP camera stream is unreachable, P1
the AI service shall POST a
CAMERA_UNAVAILABLE fault and
cease sending occupancy updates.

FR-OD-04 If YOLO and HC-SR04 classifications P2
disagree by more than one tier, the
backend shall use the higher
classification as the active value.




21
GreenLab --- Chapter 3: System Analysis


\subsection{Area: Lighting Control}

Lighting control is the first of the two automation outputs driven by the decision matrix in the
ASP.NET Core backend. These requirements specify when a lighting relay command is issued,
how an active manual override suppresses automatic lighting decisions, and how every lighting
state change is captured for traceability.

Table 3.10: Area: Lighting Control.

FR-ID Requirement Priority

FR-LC-01 The automation engine shall evaluate P1
the lighting decision matrix on every
sensor cycle and issue a relay
command only when state changes.

FR-LC-02 When an active OverrideCommand P1
exists for Relay 1 or Relay 2, the
automation engine shall not issue any
lighting relay commands.

FR-LC-03 When occupancy classification is P1
Empty and no active override exists
for light relays, the automation engine
shall command both light relays OFF.

FR-LC-04 Every relay state change shall produce P1
one row in RelayStates
(CommandSource = Automation or
Override) and one row in AuditLog.


\subsection{Area: Ventilation Control}

Ventilation control mirrors lighting control but is driven primarily by the DHT11 temperature
reading rather than occupancy alone, reflecting the fact that thermal comfort and equipment
cooling needs can apply even in an empty room. These requirements define the fan activation rule,
the priority given to the Hot temperature zone, and the same override-suppression and audit
principles applied to lighting.




22
GreenLab --- Chapter 3: System Analysis


Table 3.11: Area: Ventilation Control.

FR-ID Requirement Priority

FR-VC-01 The automation engine shall evaluate P1
the fan activation table on every sensor
cycle and command Relay 3 and Relay
4 accordingly.

FR-VC-02 When temperature zone is Hot (above P1
30$^\circ$C), the automation engine shall
command Relay 3 (Fan 1) to ON
regardless of occupancy.

FR-VC-03 When occupancy is Empty and P1
temperature zone is Cool or
Comfortable, both Relay 3 and Relay 4
shall be commanded OFF.

FR-VC-04 When an active OverrideCommand P1
exists for Relay 3 or Relay 4, the
automation engine shall not issue any
ventilation relay commands.


\subsection{Area: Monitoring and Dashboard}

These requirements define what the Angular dashboard must show, how often it refreshes, and how
visibility is scoped by role --- an instructor sees only their assigned laboratory, while an
administrator sees the full estate. They also define how fault conditions surface visually so that
they cannot be missed.

Table 3.12: Area: Monitoring and Dashboard.

FR-ID Requirement Priority

FR-MD-01 The Angular dashboard shall refresh P1
each lab state every 10 seconds via
GET /labs/{id}, showing LDR
value/zone, temperature, humidity,
occupancy, relay states.




23
GreenLab --- Chapter 3: System Analysis


FR-ID Requirement Priority

FR-MD-02 An instructor shall only view data for P1
the LabId in their JWT AssignedLabId
claim; a mismatched GET /labs/{id}
request returns HTTP 403.

FR-MD-03 An administrator shall view P2
LabStatusCards for all active
laboratories from one view, with fault
banners visible without scrolling.

FR-MD-04 Active fault events shall appear as a P1
styled banner showing fault type,
timestamp, and a Resolve button
visible only to administrators.


\subsection{Area: Manual Override}

Manual override is the mechanism through which the governance dimension of GreenLab asserts
itself over pure automation: an authorised user may temporarily take direct control of a relay.
These requirements define who may override what, how an override expires automatically, and
how administrator overrides take precedence over instructor overrides --- directly reflected in the
Manual Override sequence diagram analysed in Section 3.8.3.

Table 3.13: Area: Manual Override.

FR-ID Requirement Priority

FR-MO-01 An instructor shall issue a manual P1
override for any relay in their assigned
lab only, specifying relay target,
desired state, duration, and reason.

FR-MO-02 An administrator shall issue a manual P1
override for any relay in any lab; an
administrator override supersedes any
existing instructor override.




24
GreenLab --- Chapter 3: System Analysis


FR-ID Requirement Priority

FR-MO-03 When an override ExpiresAt is P1
reached, the backend shall set IsActive
= 0 and resume automation engine
control of the affected relay.

FR-MO-04 Every override creation, automatic P1
expiry, and manual cancellation shall
be recorded in AuditLog with ActorId,
ActorRole, and Outcome.


\subsection{Area: Fault Handling and Safety}

Because GreenLab directly switches physical appliances, the system must fail safe rather than fail
silent. These requirements define the timeout thresholds that distinguish a healthy data flow from a
stalled one, the safe state the automation engine holds during a fault, and the administrator
workflow for acknowledging and clearing a fault once resolved.

Table 3.14: Area: Fault Handling and Safety.

FR-ID Requirement Priority

FR-FH-01 The backend shall generate P1
SENSOR_TIMEOUT if no
SensorReadings row arrives within 15
seconds, and CAMERA_TIMEOUT if
no occupancy update arrives within 30
seconds.

FR-FH-02 When any fault is active, the P1
automation engine shall hold safe
state: Relay 1 ON (lights), Relays 2--4
hold last known state.

FR-FH-03 Active faults appear as fault banners P1
on the relevant LabStatusCard;
administrators see all labs; instructors
see their assigned lab only.




25
GreenLab --- Chapter 3: System Analysis


FR-ID Requirement Priority

FR-FH-04 An administrator resolves a fault via P1
PUT /faults/{id}/resolve; the backend
sets IsResolved = 1, ResolvedAt =
UTC now, ResolvedByAdminId.


\subsection{Area: Authentication and Access Control}

Every other functional area depends on this one: nothing in GreenLab is reachable without a valid
identity and a verified role. These requirements define how credentials are validated, how the
resulting JSON Web Token is structured and enforced, and the non-negotiable rule that plaintext
passwords are never stored or transmitted in clear form --- the same JWT life cycle that underlies
the authentication step of the Manual Override workflow in Section 3.8.3.

Table 3.15: Area: Authentication and Access Control.

FR-ID Requirement Priority

FR-AC-01 POST /auth/login shall validate P1
username and verify password against
bcrypt PasswordHash; on success
return a JWT signed with HMAC-
SHA256, expiring in 60 minutes.

FR-AC-02 Every endpoint except POST P1
/auth/login requires a valid non-
expired JWT bearer token;
missing/invalid token returns HTTP
401.

FR-AC-03 The backend shall enforce that an P1
instructor's JWT AssignedLabId claim
matches the labId in every request; any
mismatch returns HTTP 403.

FR-AC-04 All passwords shall be stored P1
exclusively as bcrypt hashes; plaintext
passwords shall never be stored,
logged, or returned.



26
GreenLab --- Chapter 3: System Analysis



\section{Non-Functional Requirements}

Where Section 3.6 states what the system must do, this section states how well it must do it, across
six quality dimensions: performance, reliability, security, usability, maintainability, and scalability.
Every design target below is a design target, not a measured result --- verification against these
targets belongs to the testing chapter, not this one.

\subsection{Performance}

Performance requirements bound how quickly the system must react, since a 10-second dashboard
refresh and a 5-second sensor cycle are only meaningful if every tier underneath them keeps pace.

Table 3.16: Non-functional requirements --- Performance.

NFR-ID Requirement Design Target Rationale

NFR-PERF-01 Backend responds to all API < 500 ms Dashboard 10 s polling must
requests under normal load feel instantaneous.
($\le$ 20 labs, $\le$ 50 concurrent
users).

NFR-PERF-02 End-to-end sensor-reading $\le$ 2 seconds One 5 s polling cycle must
→ relay-command latency complete within it.
under normal load.

NFR-PERF-03 Python AI service completes $\le$ 5 seconds Keeps occupancy
one full YOLOv8n classification fresh every
inference cycle per lab. polling cycle.


\subsection{Reliability}

Reliability requirements address what happens when connectivity is imperfect --- the normal
condition for IoT deployments inside a university building rather than the exception.




27
GreenLab --- Chapter 3: System Analysis


Table 3.17: Non-functional requirements --- Reliability.

NFR-ID Requirement Design Target Rationale

NFR-REL-01 ESP32 in Degraded Mode No time limit Prevents relay flapping on
maintains last known relay intermittent network.
state without reset until
connectivity returns.

NFR-REL-02 Backend uptime during lab $\ge$ 99% Faculty operations depend
operating hours (08:00-- on continuous availability.
22:00, Mon--Sat).

NFR-REL-03 On ESP32 reconnection, Within 1 poll cycle Prevents stale relay states
device retrieves current after reconnection.
relay command before
resuming automation.


\subsection{Security}

Security requirements protect the confidentiality of credentials and sensor data in transit and
guarantee that the audit trail cannot be altered after the fact.

Table 3.18: Non-functional requirements --- Security.

NFR-ID Requirement Design Target Rationale

NFR-SEC-01 All tier communication uses TLS 1.2+ Protects JWT tokens and
HTTPS; backend redirects sensor data in transit.
any HTTP request to
HTTPS.

NFR-SEC-02 Role enforcement Middleware-level Single enforcement point;
implemented at middleware no gap from missing checks.
level, not inside individual
controller methods.

NFR-SEC-03 AuditLog has no UPDATE INSERT-only Records must be tamper-
or DELETE permissions for evident and legally
any role; only INSERT by defensible.
audit middleware.




28
GreenLab --- Chapter 3: System Analysis


\subsection{Usability}

Usability requirements recognise that instructors are domain experts in teaching, not in buildingautomation software, so the override path in particular must be self-evident.

Table 3.19: Non-functional requirements --- Usability.

NFR-ID Requirement Design Target Rationale

NFR-USA-01 A first-time instructor can $\le$ 3 clicks Instructors are not IT
locate and issue a manual specialists.
override without reading
documentation.

NFR-USA-02 Fault banners visible on 1920×1080, no scroll Faults must be immediately
LabStatusCard without visible.
scrolling, up to 6 cards on
screen.


\subsection{Maintainability}

Maintainability requirements keep tuning and evolution out of the source code, consistent with the
independently-deployable four-tier philosophy described in Section 3.2.

Table 3.20: Non-functional requirements --- Maintainability.

NFR-ID Requirement Design Target Rationale

NFR-MAIN-01 All control thresholds (LDR 0 hardcoded thresholds Enables tuning without code
zones, YOLO confidence, changes or redeployment.
polling intervals, JWT
expiry, fan levels) stored in
SystemConfig.

NFR-MAIN-02 Each of the four tiers is Zero cross-tier coupling Supports iterative
independently deployable; improvement of individual
updating one tier requires no tiers.
changes to any other tier,
provided the API contract is
stable.




29
GreenLab --- Chapter 3: System Analysis


\subsection{Scalability}

Scalability requirements ensure that GreenLab can grow from a pilot deployment to the full set of
IT laboratories without a redesign.

Table 3.21: Non-functional requirements --- Scalability.

NFR-ID Requirement Design Target Rationale

NFR-SCAL-01 System supports $\ge$ 20 labs Faculty has >10 IT labs;
simultaneously active future growth must be
laboratories, each posting supported.
sensor data every 5 s,
without architecture or
schema changes.

NFR-SCAL-02 Adding a new laboratory 0 schema changes per new Reduces IT Department
requires only: one lab effort for lab additions.
Laboratories row, one user
account, one ESP32
configured, one camera
configured.




30
GreenLab --- Chapter 3: System Analysis



\section{Sequence Diagram Analysis}

A requirement table establishes correctness criteria for one behaviour in isolation. A sequence
diagram shows how several requirements, several actors, and all four tiers cooperate to deliver one
complete, time-ordered scenario. The four sequence diagrams produced for GreenLab and analysed
below were chosen because together they exercise every tier of the architecture at least once: the
occupancy workflow exercises the hardware layer, the Python AI service, and the backend; the
camera streaming workflow exercises the backend's media-handling path and the Angular
dashboard; the manual override workflow exercises authentication, the backend's automation
engine, the hardware layer, and the audit subsystem in a single pass; and the transfer request
workflow exercises the same authorisation layer applied to an organisational rather than a
hardware-level decision.

\subsection{Occupancy-Triggered Automation Workflow}

This diagram has six lifelines --- Motion Sensor, ESP32, Controller, API, Model YOLO, and
Camera --- and models the complete path from a physical presence event to an appliance being
switched, including the visual confirmation step performed by YOLO. It is the dynamic
counterpart of the Entry/Exit Detection (FR-EE), Occupancy Detection and Classification (FROD), Lighting Control (FR-LC), and Ventilation Control (FR-VC) requirement areas, and it
realises use cases UC-09, UC-10, UC-12, and UC-13.




31
GreenLab --- Chapter 3: System Analysis




Figure 3.4: Sequence diagram --- occupancy-triggered automation workflow.

The scenario begins when the Motion Sensor detects activity (step 1) and signals the ESP32, which
sends a trigger to the Controller (step 2) --- the practical instantiation of the HC-SR04 entry
classification in FR-EE-01, subject to the 3-second debounce of FR-EE-02. The Controller issues
an HTTP request to analyse the scene (step 3), which the API forwards to the Model YOLO
component as a request for frame analysis (step 4). This request-response pair is wrapped in an opt
[Camera Active] guard: if the camera is reachable, YOLO captures a frame from the Camera (step
\begin{itemize}
  \item and receives it back (step 6); if the camera is not active, this branch is skipped entirely, which is
the sequence-diagram expression of the CAMERA_UNAVAILABLE fault path defined in FROD-03 --- the diagram shows the happy path, while the fault path is specified formally by that
requirement rather than drawn as a second branch.
\end{itemize}


32
GreenLab --- Chapter 3: System Analysis


Once a frame is available, YOLO processes the image and counts students (step 7, a self-call
representing the YOLOv8n inference and the confidence-0.6 / class-0 filter of FR-OD-01), then
returns the count to the API (step 8), which relays it to the Controller as the response students
count (step 9). The Controller's next action is wrapped in an alt block with two mutually exclusive
branches that correspond directly to FR-LC-03 and FR-VC-03 on one side and FR-LC-01/FR-VC01 on the other: when [students greater than 0], the Controller tells the ESP32 to turn ON devices
(step 10), and the ESP32 activates the lights and AC/fan (step 11); when [students equal 0 for a
period], the Controller instead instructs the ESP32 to turn OFF devices (step 12), and the ESP32
deactivates the relevant components (step 13) --- this is the End of Session Reset use case (UC-12)
executing in practice. The diagram closes with a loop [Continuous Monitoring] (step 14), which is
the visual representation of the fact that this entire sequence repeats indefinitely on the 5-second
polling cycle defined across the System Initialisation and Occupancy Detection requirement areas,
rather than running once.

Two design decisions in this diagram are worth calling out explicitly because they recur in the
requirements. First, occupancy is confirmed by a vision-based count rather than acted on directly
from the motion trigger alone --- the motion event only initiates the check; the relay decision waits
for YOLO's answer. This is what allows FR-OD-04 to exist at all: a disagreement between the HCSR04-based entry/exit signal and the YOLO-based count is only meaningful because both signals
are computed and compared, not because one pre-empts the other. Second, the alt branch's OFF
path is not triggered the instant the count reaches zero but only when it is zero ``for a period'',
which is the sequence-diagram phrasing of the safe-default philosophy described in Section 3.2.3: a
single empty reading should not instantly switch off lighting that a class is merely stepping out of
for a moment.

\subsection{Live Camera Streaming Workflow}

This diagram has four lifelines --- Admin, System, CameraService, and Camera --- and models the
path an administrator follows to view a laboratory's live video feed from the Angular dashboard,
distinct from the YOLO inference path of Section 3.8.1 even though both ultimately read from the
same IP camera. It realises the visual-monitoring half of UC-02 and is underpinned, on the
implementation side, by the camera-source resolution logic described in the camera_sources.py and
droidip.py modules referenced later in the design chapter.


33
GreenLab --- Chapter 3: System Analysis




Figure 3.5: Sequence diagram --- live camera streaming workflow.

The Admin clicks ``View Camera for Lab'' (step 1), and the System requests a camera stream for
that LabId from the CameraService (step 2). This request is followed by an opt [Connection
Needed] guard: a connection to the physical Camera is established (step 3) and confirmed (step 4)
only if one is not already open, which mirrors the connection-pooling behaviour implemented by
the create_capture/_ensure_camera_started logic in the streaming service --- an already-registered
camera is reused rather than reconnected on every request. Once connected, the diagram enters a
loop [Streaming Active]: the Camera continuously sends video data to the CameraService (step 5),
which forwards the stream to the System (step 6). This loop is the sequence-diagram view of the
long-running per-camera worker thread that keeps reading frames for as long as the stream remains
registered.

The outcome is resolved by an alt block with two branches. On [Success], the System displays live
video to the Admin (step 7) --- the direct realisation of FR-MD-01's requirement that the
dashboard show live state. On [Failure], the System instead shows an error message (step 8). This
34
GreenLab --- Chapter 3: System Analysis


failure branch is the user-facing counterpart of the same condition that drives FR-OD-03 in the
occupancy workflow: an unreachable camera should produce a clear, immediate signal to the user
rather than a frozen or blank video panel, consistent with the fault-visibility requirements of FRMD-04 and FR-FH-03.

This diagram complements the occupancy workflow of Section 3.8.1 rather than duplicating it: that
diagram shows the camera being used as a sensor input to an automated decision, while this
diagram shows the same camera being used as a direct visual aid for a human administrator. Both
depend on the camera being reachable, which is why camera connectivity is treated as a fault
condition (CAMERA_UNAVAILABLE / CAMERA_TIMEOUT) in FR-OD-03 and FR-FH-01
regardless of which consumer --- YOLO or a human --- is currently using the stream.

\subsection{Manual Override (Relay Control) Workflow}

This diagram has five lifelines --- Instructor, Angular Dashboard, ASP.NET Backend, ESP32, and
Relay Module --- and is the most complete single-diagram demonstration of GreenLab's
governance dimension: in one continuous scenario it exercises authentication, role validation, relay
actuation, dashboard feedback, and audit logging. It is the dynamic realisation of use case UC-03
and of the Manual Override (FR-MO) and Authentication and Access Control (FR-AC)
requirement areas together, and it shows concretely how FR-LC-02 and FR-VC-04 (override
suppresses automatic relay commands) are enforced at the point an override is created rather than
left to be checked later.




35
GreenLab --- Chapter 3: System Analysis




Figure 3.6: Sequence diagram --- manual override (relay control) workflow.

The scenario opens with the Instructor opening the dashboard and logging in (step 1). This is
wrapped in an opt [Authentication] guard, since a returning user with an unexpired JWT can skip
straight past it: the Angular Dashboard posts the username and password to /api/auth/login (step 2),
and the ASP.NET Backend returns 200 OK with a JWT (step 3) --- the direct execution of FR-AC01, which specifies bcrypt password verification and an HMAC-SHA256-signed token expiring in
60 minutes. The dashboard then loads (step 4).

With a session established, the Instructor sets lighting to 100% through Manual Control (step 5),
and the Angular Dashboard posts this to /api/control/manual-override with the JWT attached as a
Bearer token (step 6). The backend's first action is to validate the JWT and confirm the caller's
role (step 7, a self-call) --- this is FR-AC-02 and FR-AC-03 being enforced at the middleware level
as required by NFR-SEC-02, before any relay logic runs, and it is also the point at which FR-MO01's constraint that an instructor may only override a relay in their own assigned lab would be
checked. Only once that check passes does a loop [Device Command] begin: the backend sends the
payload (lighting 100, mode MANUAL) to the ESP32 (step 8), the ESP32 performs the relay
36
GreenLab --- Chapter 3: System Analysis


actuation --- channel ON, intensity 100 (step 9) --- against the physical Relay Module, receives an
acknowledgement (step 10), and forwards an ACK that the command was executed back to the
backend (step 11). This loop framing reflects that a real override may legitimately need to retry or
re-issue the device command, not that the user repeats their request.

Once the device confirms execution, the backend updates its dashboard data store (step 12, a selfcall corresponding to writing a new RelayStates row with CommandSource = Override, satisfying
FR-LC-04) and returns 200 OK with state MANUAL and lighting 100 to the Angular Dashboard
(step 13). The outcome is then resolved by an alt block: on [Success], the dashboard displays a
MANUAL badge and updates its indicators (step 14), giving the instructor immediate visual
confirmation that automation has been suspended for that relay, exactly as FR-LC-02 requires; on
[Failure], the dashboard shows an error message instead (step 15), consistent with the same failvisibly principle used in the camera streaming workflow of Section 3.8.2. Independently of which
branch was taken, the backend logs the action to the audit table with timestamp, user ID, action,
and resulting state (step 16, a self-call) --- the direct execution of FR-MO-04 and FR-LC-04, and
the reason the AuditLog table is defined as INSERT-only under NFR-SEC-03: this final step is not
optional or conditional on success, because both a successful and a failed override attempt are
governance-relevant events that must be attributable.

Taken together with the other diagrams, this workflow demonstrates the precedence rule of FRMO-02 implicitly: because the override-creation path runs through the same role-validation selfcall at step 7, an administrator's request would be authorised to target any laboratory's relay, while
an instructor's request is constrained to their own AssignedLabId, and a later administrator
override on the same relay would simply repeat steps 6--16 and overwrite the active
OverrideCommand row, which is why FR-MO-02 states that an administrator override supersedes
an existing instructor override rather than being rejected by it.

\subsection{Laboratory Transfer Request Workflow}

This diagram has five lifelines --- Instructor, Angular Dashboard, ASP.NET Backend,
Administrator, and Notification Service --- and models the governance workflow through which an
instructor formally requests that their session be moved to a different laboratory, and an
administrator reviews and resolves that request. It is the dynamic realisation of use case UC-06
(Approve/Reject Transfer Requests) together with UC-04 (Submit Laboratory Transfer Request),

37
GreenLab --- Chapter 3: System Analysis


and it exercises the TransferRequests table, the role-based authorisation layer (FR-AC), and the
audit logging requirement (FR-MO-04) in a single end-to-end scenario.




Figure 3.7: Sequence diagram --- laboratory transfer request workflow.

The scenario opens with the Instructor logging in and opening the dashboard (step 1), protected by
the same opt [Authentication] guard used in Section 3.8.3: the Angular Dashboard posts credentials
to /api/auth/login (step 2), the ASP.NET Backend returns 200 OK with a JWT (step 3), and the
dashboard loads (step 4). This authentication preamble is identical across all instructor-initiated
workflows because FR-AC-01 applies uniformly --- no action is reachable without a verified
identity and a role-scoped token.

With a session established, the Instructor submits a transfer request through the dashboard form
(step 5), and the Angular Dashboard posts it to /api/transfer/request with the JWT as a Bearer token
(step 6). The backend validates the JWT and confirms the instructor role (step 7, a self-call) ---
enforcing FR-AC-02 and FR-AC-03 at the middleware level before any business logic runs. Once
authorised, the backend creates a new TransferRequests row with status PENDING (step 8, a selfcall), and returns 201 Created to the Angular Dashboard (step 9), which confirms the submission to
38
GreenLab --- Chapter 3: System Analysis


the Instructor (step 10). This creation step is the direct execution of the Submit Laboratory Transfer
Request use case (UC-04).

The workflow then shifts to the administrator's side, modelled here as a separate lifeline that picks
up asynchronously. A loop [Pending Requests] reflects that an administrator may review one or
many pending requests in a single session. The Notification Service (step 11) alerts the
Administrator that a new transfer request is awaiting review, and the Administrator retrieves the
request details from the backend via GET /api/transfer/{id} (step 12), receiving the full
TransferRequests record in response (step 13). The administrator then makes a resolution decision,
captured by an alt block: on [Approve], the Angular Dashboard puts /api/transfer/{id}/approve
(step 14), the backend updates the TransferRequests row to APPROVED (step 15, a self-call), and
the Notification Service sends an approval notification to the Instructor (step 16); on [Reject], the
Angular Dashboard puts /api/transfer/{id}/reject (step 17), the backend updates the row to
REJECTED (step 18, a self-call), and the Notification Service sends a rejection notification instead
(step 19). This branching realises the Approve/Reject Transfer Requests use case (UC-06) and its
<<extend>> relationship: approval and rejection are both outcomes of the same review action rather
than separate triggers.

Independently of the resolution branch, the backend logs the entire transaction --- request creation,
review, and outcome --- to the AuditLog table with timestamp, userID, action, and resulting state
(step 20, a self-call), satisfying FR-MO-04 and the INSERT-only constraint of NFR-SEC-03. As in
the manual override workflow, this logging step is unconditional: a rejected transfer request is as
governance-relevant as an approved one, because the audit trail must be able to show not only what
changes were made to laboratory assignments, but also which requests were considered and
declined.

Together, the four sequence diagrams of Section 3.8 demonstrate that GreenLab's governance
model --- role-scoped tokens, explicit decision branches, mandatory audit records --- applies
uniformly across automated, administrative, and instructor-initiated workflows rather than being
confined to any single scenario.




39
GreenLab --- Chapter 3: System Analysis



\section{Chapter Summary}

This chapter established the analytical basis for the GreenLab system. Section 3.2 positioned
GreenLab as a four-tier cyber-physical system shaped by two forces beyond simple sensor-to-relay
automation: a temporal dimension driven by irregular academic scheduling, and a governance
dimension driven by institutional accountability requirements --- and showed how both are
answered by safe defaults at every tier rather than by assuming reliable connectivity. Section 3.3
then positioned that same architecture commercially, using the Business Model Canvas to show
that GreenLab's automation behaviour, governance controls, and reliability targets double as the
value propositions, cost drivers, and revenue streams of a viable product rather than a prototype
alone.

Section 3.4 identified the six stakeholder groups --- Instructor, Administrator, Laboratory Staff,
Student, IT Department, and Faculty Management --- whose concerns motivate every feature in the
system. Section 3.5 converted those concerns into a formal use case model of thirteen use cases
across two human and two automated actors, distinguishing governance actions that require an
authenticated session from automation actions that run continuously.

Sections 3.6 and 3.7 then fixed the testable requirement baseline for the rest of the project: thirtysix functional requirements across the nine areas of System Initialisation and Startup, Entry and
Exit Detection, Occupancy Detection and Classification, Lighting Control, Ventilation Control,
Monitoring and Dashboard, Manual Override, Fault Handling and Safety, and Authentication and
Access Control; and the non-functional requirements spanning performance, reliability, security,
usability, maintainability, and scalability. Every requirement identifier, threshold, and design target
stated in these two sections is taken verbatim from the TA Reference Document and is treated in
this project as the single source of truth, taking precedence over any earlier draft figures that may
appear elsewhere.

Finally, Section 3.8 validated that requirement baseline dynamically by walking through the four
sequence diagrams produced for GreenLab. The occupancy-triggered automation workflow
showed how a motion event, a YOLO-confirmed person count, and a debounce-aware decision
combine to switch lighting and ventilation automatically. The live camera streaming workflow
showed how the same camera infrastructure also serves a human administrator directly, with an
explicit failure path rather than a silent one. The manual override workflow tied authentication, role
40
GreenLab --- Chapter 3: System Analysis


validation, relay actuation, dashboard feedback, and audit logging into a single traceable scenario,
demonstrating concretely how an override suppresses automation and how administrator
precedence over instructor overrides is structurally guaranteed rather than merely documented. The
laboratory transfer request workflow completed the picture by modelling the institutional
governance process through which an instructor requests a room change and an administrator
approves or rejects it, demonstrating that the same role-scoped authorisation, conditional
branching, and unconditional audit logging that govern hardware-level actions apply equally to
organisational workflows.

Together, the business model, the stakeholder and use case analysis, the requirement tables, and the
sequence diagram analysis in this chapter form the contract that the design chapter must satisfy:
every component introduced in Chapter 4 --- the database schema, the REST API, the automation
engine, the AI service, and the Angular dashboard --- exists because a specific requirement, use
case, or sequence-diagram step in this chapter demands it.




41
